\scanpage{105}
\begin{center}\uline{Fidélité}\end{center}

% slanted
\margin{(ancienne version)}
\deleted{\uline{\unsure{Corollaire}}} \inserted{Proposition} Soit une réalisation géométrique
régulière $\varphi$ \inserted{(non dégénérée (et propre))} d'un polygone combinatoire
$\Pi\inserted{=(\underline S,\underline A,\underline R)}$ dans un plan affine $E$ \inserted{sur un corps $k$}. Conditions équiv.:

a) $\varphi$ est fidèle

b) $\varphi_S:\underline S\to E$ injective

b$'$) $\varphi_A:\underline A\to\mathrm{Dr}(E)$ injective

c) $\mathrm{Aut}(\Pi)=G\unsure{-|\to}\mathrm{Aff}(E)$ injective

d) \deleted{[Si $\Pi$ orienté, d'où $\underline u\in\mathrm{Aut}(\Pi)$, et $u\in\mathrm{Aff}(E)$]}
\deleted{$\mathrm{Aut}^{+}(\Pi)$} $G^+\unsure{-|\to}\mathrm{Aff}(E)$ injective

\uline{dém} a) $\Rightarrow$ b) \& b$'$) trivial (\unsure{compte|cpte} tenu que
$\forall\underline s,\underline s'\in\underline S$, $\underline s\ne\underline s'$, $\exists\underline a\in\underline A$ qui est incident à $\underline s$ et
non à $\underline s'$ -- car $\mathrm{card}\,S=\mathrm{card}\,A\geq3$ ici ! -- \inserted{et dualement.}

b) $\Rightarrow$ c), \inserted{b$'$) $\Rightarrow$ c)} \unsure{triviales|trivial|triviaux}, puisque $G\to\unsure{\mathfrak S|\mathcal S}_S$, \inserted{$G\to\mathfrak S_A$} injectives.

\deleted{\unsure{Prouvons}} c) $\Rightarrow$ d) trivial. Prouvons
d) $\Rightarrow$ b), i.e. non b $\Rightarrow$ non d). Choisissons
\deleted{un gén.} repère $r$ de $\Pi$, d'où un gén.
\deleted{\unsure{u de}} $\underline u=\underline\sigma_0\underline\sigma_1$ de $\underline G=\mathrm{Aut}(\Pi)$, d'où $u\in\mathrm{Aff}(E)$.
\deleted{On aura} Soit $n=\mathrm{Card}\,\underline S$ (\deleted{de sorte que \unsure{les}} \inserted{$3\leq n\leq+\infty$})
\deleted{Dire que} \inserted{Alors} non b signifie $\exists d<n$ \deleted{tel que}
\deleted{$\underline s_d$} et $i\in\mathbb Z$ tels que $s_i=s_{i+d}$ i.e.
$u^is_0=u^{i+d}s_0$ i.e. \deleted{\unsure{$s$}} $u^ds_0=s_0$ ce qui
implique, en multipliant par $u^j$ et
notant $u^ju^ds_0\unsure{-|=}u^du^js_0=u^ds_j$, que
$u^ds_j=s_j$ $\forall j\in\mathbb Z$, et comme les $s_j$ engendrent
$E$, on aura $u^{d}=\mathrm{id}$, \deleted{\unsure{donc}} or $\underline u^d\ne1$, donc
non d). On a ainsi prouvé
% archive sheet number: 104
